\reading{CR-15}{Counterparty Risk and Beyond}{DEFER}{0}

\begin{crkeybox}[Learning objectives for this reading]
\begin{enumerate}[label=\textbf{\alph*.},leftmargin=2em]
\item Describe counterparty risk and differentiate it from lending risk.
\item Describe transactions that carry counterparty risk and explain how counterparty risk can arise in each transaction.
\item Identify and describe institutions that take on significant counterparty risk.
\item Describe credit exposure, credit migration, recovery, mark-to-market, replacement cost, default probability, loss given default, and the recovery rate.
\item Describe credit value adjustment (CVA) and compare the use of CVA and credit limits in evaluating and mitigating counterparty risk.
\item Identify and describe the different ways institutions can quantify, manage, and mitigate counterparty risk.
\item Identify and explain the costs of an OTC derivative.
\item Explain the components of the X-Value Adjustment (xVA) term.
\end{enumerate}
\end{crkeybox}

% =====================================================================
\lo{CR-15 a}{Describe counterparty risk and differentiate it from lending risk.}

\begin{center}
\footnotesize
\begin{tabular}{L{2.7cm} L{5.8cm} L{6.0cm}}
\toprule
 & \thead{Counterparty risk} & \thead{Lending risk} \\
\midrule
Definition & The risk that the other party to a financial contract will fail to
  fulfil their side of the agreement &
  The risk that the other party fails to pay some or all of a loan, due to
  insolvency \\
Where it exists & OTC derivatives (interest rate swaps, FX forwards, CDS) and
  securities financing transactions (repos, reverse repos, securities borrowing
  and lending) &
  Loans, bonds, mortgages, credit cards \\
Amount at risk & \textbf{Not constant}; the exposure varies between positive and
  negative &
  The notional amount at risk is usually known with a degree of certainty \\
Who bears it & \textbf{Bilateral}: both parties take risk &
  \textbf{Unilateral}: only one party takes risk, and it primarily affects the
  lender \\
\bottomrule
\end{tabular}
\end{center}
\figcap{The two rows that decide most questions are the bottom two. Lending risk
is one-directional with a known principal; counterparty risk is two-directional
with an exposure that swings with the value of the underlying.}

\begin{crdefbox}[Pre-settlement and settlement risk]
\term{Pre-settlement risk} is the risk that a counterparty will default
\emph{prior} to the final settlement of the transaction, at expiration. This is
what counterparty risk usually refers to.

\smallskip
\term{Settlement risk} arises \emph{at} final settlement, if there are timing
differences between when each party performs on its obligations under the
contract.
\end{crdefbox}

% =====================================================================
\lo{CR-15 b}{Describe transactions that carry counterparty risk and explain how counterparty risk can arise in each transaction.}

\subsection*{OTC derivatives}

Because of customization --- to reduce basis risk --- a much greater amount of
notional is traded OTC. Interest rate swaps, FX forwards and credit default swaps
all carry counterparty risk.

\begin{itemize}
\item \term{Interest rate swaps.} Risk arises from the exchange of floating
versus fixed cash flows.
\item \term{Foreign exchange transactions.} Risk arises from the large notional
amounts exchanged.
\item \term{Credit default swaps.} Introduce additional risk through
\term{wrong-way risk}, where a credit rating downgrade increases the likelihood
of the counterparty defaulting at the same time as the protection is needed.
\end{itemize}

\subsection*{Securities financing transactions}

Repos, reverse repos, and securities borrowing and lending. Risk exists if the
borrower fails to repurchase the securities, or if the collateral value declines.

\begin{center}
\begin{tikzpicture}[font=\sffamily\footnotesize]
% leg 1
\node[font=\scriptsize\bfseries, text=FRMNavy] at (2.6,2.35) {Opening leg};
\node[draw=FRMNavy, fill=PaleNavy, line width=0.9pt, circle, minimum size=0.95cm] (a1) at (0,1.1) {\textbf{A}};
\node[draw=FRMNavy, fill=PaleNavy, line width=0.9pt, circle, minimum size=0.95cm] (b1) at (5.2,1.1) {\textbf{B}};
\draw[-{Stealth[length=5pt]}, FRMGold, line width=1pt] ($(a1.east)+(0,0.28)$) --
  node[above, font=\scriptsize, inner sep=1.5pt]{Gov sec \$100M} ($(b1.west)+(0,0.28)$);
\draw[-{Stealth[length=5pt]}, FRMGold, line width=1pt] ($(b1.west)-(0,0.28)$) --
  node[below, font=\scriptsize, inner sep=1.5pt]{cash \$100M} ($(a1.east)-(0,0.28)$);
% leg 2
\node[font=\scriptsize\bfseries, text=FRMNavy] at (10.4,2.35) {Closing leg, at maturity};
\node[draw=FRMNavy, fill=PaleNavy, line width=0.9pt, circle, minimum size=0.95cm] (a2) at (7.8,1.1) {\textbf{A}};
\node[draw=FRMNavy, fill=PaleNavy, line width=0.9pt, circle, minimum size=0.95cm] (b2) at (13.0,1.1) {\textbf{B}};
\draw[-{Stealth[length=5pt]}, FRMGreen, line width=1pt] ($(b2.west)+(0,0.28)$) --
  node[above, font=\scriptsize, inner sep=1.5pt]{Gov sec \$100M} ($(a2.east)+(0,0.28)$);
\draw[-{Stealth[length=5pt]}, FRMGreen, line width=1pt] ($(a2.east)-(0,0.28)$) --
  node[below, font=\scriptsize, inner sep=1.5pt]{cash \$100M $+$ \$1M interest} ($(b2.west)-(0,0.28)$);
\end{tikzpicture}
\end{center}
\figcap{A repo, both legs. A sells securities for cash and buys them back later at
a higher price; the difference is the repo interest. The counterparty risk is not
in either leg on its own but in the gap between them: B is exposed if A fails to
repurchase, and both are exposed to a move in the collateral value in the
meantime.}

\begin{crkeybox}[The exclusion worth remembering]
Exchange transactions do not carry counterparty risk.
\end{crkeybox}

% =====================================================================
\lo{CR-15 c}{Identify and describe institutions that take on significant counterparty risk.}

\begin{center}
\footnotesize
\begin{tabular}{L{2.5cm} L{5.9cm} L{6.1cm}}
\toprule
\thead{Tier} & \thead{Who they are} & \thead{Characteristics} \\
\midrule
\term{Large derivatives player} & Typically a large bank, often known as a
  \emph{dealer} &
  A vast number of OTC derivatives trades on their books; trade with each other
  and have many other clients; coverage of all or many asset classes; post
  collateral against positions \\
\term{Medium derivatives player} & Typically a smaller bank or other financial
  institution, such as a hedge fund or pension fund &
  Many OTC derivatives trades; trade with a relatively large number of clients;
  cover several asset classes though may not be active in all; probably post
  collateral, with some exceptions \\
\term{Small derivatives player} & A large corporate or sovereign with significant
  derivatives requirements &
  A few OTC derivatives trades; potentially only a few counterparties; may be
  specialized to a single asset class; \textbf{unable to commit to posting
  collateral, or will post illiquid collateral} \\
\bottomrule
\end{tabular}
\end{center}
\figcap{The collateral column is the discriminating one, and it inverts the
intuition: the smallest players, with the least capacity to absorb a loss, are
also the ones least able to post collateral against it.}

% =====================================================================
\lo{CR-15 d}{Describe credit exposure, credit migration, recovery, mark-to-market, replacement cost, default probability, loss given default, and the recovery rate.}

\begin{enumerate}
\item \term{Credit exposure.} The potential loss if the counterparty defaults,
considering current, future and contingent liabilities.
\item \term{Credit migration.} A counterparty's creditworthiness can change over
time. \emph{Deterioration} means increased default risk, especially for long-term
contracts. \emph{Improvement} means reduced default risk, and is more likely for
initially weak credit ratings.
\item \term{Recovery rate.} The portion of the outstanding debt recoverable after
a default. A 70\% recovery rate means a 30\% loss.
\item \term{Loss given default (LGD).} One minus the recovery rate:
$\mathrm{LGD} = 1 - \mathrm{RR}$.
\item \term{Mark-to-market (MtM).} The current potential loss on all contracts
with a counterparty, considering present values of future cash flows. Not a
perfect measure, because it ignores netting, collateral and hedging.
\item \term{Replacement cost.} The cost to replace the defaulted contract in the
current market.
\item \term{Default probability.} Treated in full at CR-9, where it is estimated
from ratings, hazard rates, credit spreads and the Merton model.
\end{enumerate}

\begin{crtrapbox}[Credit migration cuts both ways]
Improvement is more likely for \emph{initially weak} ratings and deterioration is
more likely over \emph{long-term} contracts. Both directions are examinable and
the conditions attached to each are swappable.
\end{crtrapbox}

% =====================================================================
\lo{CR-15 e}{Describe credit value adjustment (CVA) and compare the use of CVA and credit limits in evaluating and mitigating counterparty risk.}

CVA is the valuation of counterparty risk on a bilateral trade: the present value
of the expected cost to the institution of a counterparty default. The full
definition and the credit-adjusted value identity are at \textbf{CR-12 b}.

\begin{crkeybox}[CVA versus credit limits: they answer different questions]
\term{CVA} helps you pick good counterparties and optimize netting: it prices the
risk.

\smallskip
\term{Credit limits} ensure you do not overexpose yourself to any single
counterparty: they cap the risk.
\end{crkeybox}

\subsection*{Quantifying counterparty risk, by level}

\begin{enumerate}[label=\alph*)]
\item \term{Trade level.} Analyse the specific trade and its risk factors.
\item \term{Counterparty level.} Consider risk-mitigating factors like netting and
collateral for each counterparty.
\item \term{Portfolio level.} Evaluate overall exposure, assuming limited defaults
within a set period.
\end{enumerate}

% =====================================================================
\lo{CR-15 f}{Identify and describe the different ways institutions can quantify, manage, and mitigate counterparty risk.}

\subsection*{Managing counterparty risk}

\begin{enumerate}[label=\alph*)]
\item \term{Choose high-quality counterparties.} The simplest method. Look for
counterparties with strong credit ratings, who may not even need to provide
collateral.
\item \term{Cross-product netting.} Aggregate derivative transactions to net out
exposures, reducing overall risk.
\item \term{Close-out.} Immediately cancel all contracts with a defaulting
counterparty, then net the values to determine who owes whom.
\item \term{Collateralization.} Require posting of collateral to cover net
exposures, \emph{theoretically} reducing exposure to zero.
\item \term{Walkaway features.} Clauses allowing cancellation of a transaction if
the other party defaults, especially where that benefits you financially.
\item \term{Diversification.} Spread counterparty risk across a variety of
institutions, reducing exposure to any single default.
\item \term{Centralized clearinghouses.} Use exchanges that act as counterparty to
all trades, centralizing and potentially reducing risk.
\end{enumerate}

\subsection*{Mitigating counterparty risk}

\begin{enumerate}
\item \term{Netting.} Offsetting negative and positive contract values with the
same counterparty, so only the net debtor pays. Reduces the amount outstanding at
any one time.
\item \term{Collateral.} A deposit of cash or securities held against the value of
the transaction. Theoretically eliminates risk, but carries administrative and
liquidity risks.
\item \term{Hedging.} Employ credit derivatives to shift risk, \emph{at the cost
of creating market risk}. Less likely in practice for an end user: if a company is
facing an investment bank on a rate swap it presumably has a legitimate business
use for the swap. The bank can hedge that risk with another counterparty, but risk
facing the smaller company is not eliminated.
\item \term{Central counterparties (CCPs).} Centralize risk via clearinghouses. In
theory this eliminates counterparty risk, but in a large systemic failure the CCP
itself could be brought down, and it may create operational and systemic risks.
\item \term{Other clauses.} Features such as resets or additional termination
events, which periodically reset MtM values or terminate transactions early.
\end{enumerate}

% =====================================================================
\lo{CR-15 g}{Identify and explain the costs of an OTC derivative.}

The costs depend on which side of the trade you are on.

\begin{center}
\footnotesize
\begin{tabular}{L{3.0cm} L{11.5cm}}
\toprule
\thead{Positive MtM} & Your trade is profitable, and you are in the money. \\
 & \term{Counterparty risk}: you are exposed if the counterparty defaults,
especially on the \emph{uncollateralized} portion. \\
 & \term{Funding cost}: you might incur funding costs for the collateralized
part, depending on who chooses the collateral. \\
\midrule
\thead{Negative MtM} & Your trade is losing money, and you are out of the money. \\
 & \term{Counterparty risk}: \emph{you} could default, if your own financial
situation worsens. \\
 & \term{Funding benefit}: you earn interest on any collateral you do not post.
You might still incur funding costs on the collateralized part, depending on who
chooses the collateral. \\
\bottomrule
\end{tabular}
\end{center}
\figcap{The symmetry is the point. When you are in the money the counterparty risk
is \emph{theirs to create and yours to bear}; when you are out of the money it is
yours to create. Funding runs the opposite way: a cost when in the money, a
benefit when out.}

% =====================================================================
\lo{CR-15 h}{Explain the components of the X-Value Adjustment (xVA) term.}

xVA is a family of terms representing the various costs and benefits associated
with OTC derivatives.

\begin{center}
\footnotesize
\begin{tabular}{L{1.8cm} L{8.6cm} L{3.6cm}}
\toprule
\thead{Term} & \thead{What it is} & \thead{Sign} \\
\midrule
\term{CVA} & Credit value adjustment: the valuation of counterparty risk on a
  bilateral trade; the cost of the \emph{other party's} default &
  \textbf{Negative} (a cost) \\
\term{DVA} & Debt value adjustment: the counterpart to CVA, reflecting
  \emph{your own} potential default risk &
  \textbf{Positive} (a benefit) \\
\term{FVA} & Funding value adjustment: the cost or benefit arising from funding
  the transaction. A trade that partially hedges a book may offer funding relief,
  since the net risk of the book is reduced &
  Either, depending on the trade \\
\term{ColVA} & Collateral value adjustment: the cost or benefit resulting from
  embedded options in collateral agreements, and any additional collateral terms &
  Either, depending on the trade \\
\term{KVA} & Capital value adjustment: the cost of holding regulatory capital over
  the entire life of the trade &
  \textbf{Always a cost} \\
\term{MVA} & Margin value adjustment: the cost of posting initial margin over the
  life of the trade &
  \textbf{Always a cost} \\
\bottomrule
\end{tabular}
\end{center}
\figcap{More than the definitions themselves, the examinable content is the sign
column: which terms are a debit and which a credit. Note that KVA and MVA are
named for capital and \emph{initial} margin specifically, and neither can ever be
a benefit.}

\begin{crtrapbox}[The sign that is most often flipped]
CVA is \term{negative} --- it is the cost of the counterparty defaulting on you.
DVA, which is \emph{your own} default, is \term{positive}. FVA and ColVA can be
either, depending on the trade. The costs associated with regulatory capital (KVA)
and initial margin (MVA) are always going to be costs.
\end{crtrapbox}
